% !PRINTED 51
\weberpage{51}
%<<CONT

\begin{equation}
 (x+y+z+\cdots)^n=\Pi(n)\sum^{\alpha,\nu}{x^\alpha u^\nu
\over \Pi(\alpha)\Pi(\nu)},
\tag{6}
\end{equation}
with the restriction
\begin{equation}
 \alpha+\nu=n.
\tag{7}
\end{equation}
By the assumption, however, it has already been proved that
\begin{equation}
 u^\nu=\Pi(\nu)\sum^{\beta,\gamma,\ldots}
 {y^\beta z^\gamma\cdots
\over \Pi(\beta)\Pi(\gamma)\cdots},
\tag{8}
\end{equation}
where
\begin{equation}
 \beta+\gamma+\cdots=\nu.
\tag{9}
\end{equation}
When this is inserted in (6), formula (3) is obtained immediately, and (7) passes over into (4).

The coefficients
\begin{equation}
 P_{\alpha,\beta,\gamma,\ldots}^{(n)}={\Pi(n)\over
 \Pi(\alpha)\Pi(\beta)\Pi(\gamma)\cdots},
\tag{10}
\end{equation}
which, according to their meaning, are whole numbers, are called the \gesperrt{polynomial coefficients}.

For example, for the third power of the trinomial one obtains
\begin{equation}
\begin{aligned}
(x+y+z)^3={}&x^3+y^3+z^3+3x^2y+3xy^2+3x^2z+3xz^2\\
&+3y^2z+3yz^2+6xyz.
\end{aligned}
\tag{11}
\end{equation}


\begin{center}
{\large \S.\ 13.}

\vspace{0.5em}
{\large\gesperrt{Derived Functions.}}
\end{center}

Let
\begin{equation}
 f(x)=a_0x^n+a_1x^{n-1}+a_2x^{n-2}+\cdots+a_n
\tag{1}
\end{equation}
be a whole rational function of order $n$.

If in it we replace $x$ by a binomial $x+y$, then we can apply the binomial theorem to each separate term, and can arrange the result according to descending or according to ascending powers of $x$ or of $y$. We shall carry out the arrangement according to ascending powers of $y$. The highest power of $y$ which occurs is the $n$th, and the coefficient of the zeroth power of $y$ is the function $f(x)$ itself, as one sees by putting $y=0$. We therefore put, denoting the other coefficients by
\[
 f'(x),\qquad {f''(x)\over \Pi(2)},\qquad {f'''(x)\over \Pi(3)},\ldots:
\]
%CONT>>
